\par\addvspace{7pt}\Needspace{4\baselineskip}\noindent\textbf{课堂题：树、有向图与调度}\par
\begin{exercise}\label{mit:6042-cp10-q4}
\textbf{6.042J，Spring 2015，CP10 第 4 题（补充题）。}
叶带标签的满二叉树递归定义为 $\langle\mathrm{leaf},\ell\rangle$，或 $\langle\mathrm{bintree},G_1,G_2\rangle$。大小满足叶的大小为 $1$，构造树的大小为 $|G_1|+|G_2|+1$。原图有四叶，从左至右为 win、lose、win、win：根的右子树是最后一个叶，左子树的左孩子是第一个叶，右孩子是中间两叶组成的二叉树。
\begin{enumerate}[label=(\alph*)]
\item 用上述括号记号写出这棵树。
\item 递归定义 $\mathrm{flatten}(G)$ 为从左到右的叶标签序列。
\item 用结构归纳证明 $2\,\mathrm{length}(\mathrm{flatten}(G))=|G|+1$。
\end{enumerate}
\end{exercise}
\sourceof{CP10，物理第 3--4 页，图 1；编者独立解答。}
\begin{solution}
(a) 树为
\[
\left\langle\mathrm{bintree},
 \left\langle\mathrm{bintree},\langle\mathrm{leaf},\mathrm{win}\rangle,
  \langle\mathrm{bintree},\langle\mathrm{leaf},\mathrm{lose}\rangle,
  \langle\mathrm{leaf},\mathrm{win}\rangle\rangle\right\rangle,
 \langle\mathrm{leaf},\mathrm{win}\rangle\right\rangle.
\]
(b) 叶的展开为单项序列 $(\ell)$；构造树的展开为 $\mathrm{flatten}(G_1)$ 与 $\mathrm{flatten}(G_2)$ 的拼接。

(c) 叶情形为 $2\cdot1=1+1$。构造步由归纳假设得到
$2\,\mathrm{length}(\mathrm{flatten}(G))=(|G_1|+1)+(|G_2|+1)=|G|+1$。
标签可以重复，因为这里数的是叶的位置数，不是不同标签数；这与 PS4 第 2 题的计数对象不同。
\end{solution}

\begin{exercise}\label{mit:6042-cp16-q1}
\textbf{6.042J，Spring 2015，CP16 第 1 题。}
\begin{enumerate}[label=(\alph*)]
\item 构造一个顶点 $v$ 在正偶长度闭合游走上、但全图没有偶长度有向圈的例子。
\item 构造一个 $v$ 在奇长度闭合游走上、却不在奇长度有向圈上的例子。
\item 证明每个奇长度闭合游走中至少有一个顶点在奇长度有向圈上。
\end{enumerate}
\end{exercise}
\sourceof{CP16，物理第 1 页；编者独立解答。}
\begin{solution}
(a) 定向三圈 $v\to a\to b\to v$ 走两遍得到长 $6$ 的闭合游走，图的简单有向圈只有长 $3$ 的圈。

(b) 弧为 $v\to c,c\to v,c\to a,a\to b,b\to c$。游走 $v,c,a,b,c,v$ 长 $5$；经过 $v$ 的简单圈却只能是 $v,c,v$，长 $2$。

(c) 对游走长度归纳。若除首尾外没有重复顶点，它本身就是奇圈。否则选一个内部重复顶点，把游走拆成两个较短的闭合游走；两个长度之和为奇数，故其中一个长度为奇数。对这个较短奇闭合游走用归纳假设，得到所需奇圈，其顶点也是原游走顶点。
\end{solution}

\begin{exercise}\label{mit:6042-cp16-q2}
\textbf{6.042J，Spring 2015，CP16 第 2 题。}
有向图中设 $u$ 可达 $v$，距离按弧数计算。证明
$\mathrm{dist}(u,v)=\mathrm{dist}(u,x)+\mathrm{dist}(x,v)$ 当且仅当 $x$ 位于某条 $u$ 到 $v$ 的最短路径上：（a）证明正向；（b）证明反向。原题的有限距离条件在这里明确写出，以排除 $\infty=\infty$ 却没有路径的情形。
\end{exercise}
\sourceof{CP16，物理第 1 页；编者独立解答，补明可达条件。}
\begin{solution}
(a) 等式左端有限，故右端两项有限。拼接一条 $u$ 到 $x$ 最短路径和一条 $x$ 到 $v$ 最短路径，得到长度等于 $\mathrm{dist}(u,v)$ 的游走。若有重复顶点，可删去正长度闭合段而得到更短游走，再删重复可成路径，违反最短距离定义。因此拼接已是经过 $x$ 的最短路径。

(b) 将最短路径在 $x$ 处分成两段。任一段若不是相应端点间的最短路径，替换成更短路径后会得到更短 $u$ 到 $v$ 游走，再删重复成路径，矛盾。因此两段长度分别为右端的两项，其和为左端。
\end{solution}

\begin{exercise}\label{mit:6042-cp16-q3}
\textbf{6.042J，Spring 2015，CP16 第 3 题。}
称包含全部八种连续三位二进制串的字符串为 $3$-好串。
\begin{enumerate}[label=(\alph*)]
\item 找一个长度 $10$ 的 $3$-好串，并证明最短长度为 $10$。
\item 在顶点 $00,01,10,11$ 的有向图上，从 $ab$ 有弧到 $b0$ 和 $b1$，弧标新增位。说明覆盖全部弧的游走如何给出好串，并给出 (a) 的对应游走。
\item 说明每条弧恰用一次的游走为何给出最短好串。
\item 推广到 $k\ge2$：顶点是 $k$ 位串，定义 $B_k$ 的弧，求最短 $(k+1)$-好串长度，证明每点入度和出度为偶数，任意有序顶点对（含同一点）之间有长至多 $k$ 的正长度路径或闭合游走。
\end{enumerate}
\end{exercise}
\sourceof{CP16，物理第 1--2 页；编者独立解答。(d) 同一点之间按正闭合游走理解。}
\begin{solution}
(a) 可取 $0001011100$。八个窗口依次为 $000,001,010,101,011,111,110,100$。长度 $L$ 的串只有 $L-2$ 个三位窗口，故 $L-2\ge8$，即 $L\ge10$。

(b) 先写游走起点的两位，再依次写每条弧的新增位。每条弧 $ab\to bc$ 对应三位窗口 $abc$。上述串对应
$00,00,01,10,01,11,11,10,00$。

(c) 图有 $8$ 条弧；每弧恰一次给 $8+2=10$ 位，并覆盖八种窗口，所以达到 (a) 下界。

(d) 弧为 $a_1\cdots a_k\to a_2\cdots a_kb$，$b\in\{0,1\}$。每点入度、出度均为 $2$，共有 $2^{k+1}$ 条弧。依次追加目标串的 $k$ 位能在 $k$ 步到达目标，故强连通；同点也有这样的正闭合游走，不同点的游走删除闭合段可成不长于 $k$ 的路径。由定理~\ref{thm:directedeuler} 存在欧拉闭合游走。将其展开得到长 $2^{k+1}+k$ 的好串；任何长 $L$ 的串至多有 $L-k$ 个窗口，故这正是最短长度。源题对同一点称“positive path”，本册采用路径无重复顶点的约定，故须作上述区分。
\end{solution}

\begin{exercise}\label{mit:6042-cp16-q4}
\textbf{6.042J，Spring 2015，CP16 第 4 题（补充题）。}
竞赛图的“排名”是经过每个顶点一次的有向路径。（a）举有多个排名的竞赛图；（b）证明每个有限竞赛图都有排名。
\end{exercise}
\sourceof{CP16，物理第 2 页；编者独立解答。}
\begin{solution}
(a) 定向三圈 $a\to b\to c\to a$ 有 $a,b,c$；$b,c,a$；$c,a,b$ 三种排名。

(b) 对顶点数归纳。删去 $x$，由归纳得到 $v_1\to\cdots\to v_{n-1}$。若 $x\to v_1$，将 $x$ 放在首位；若每个 $v_i\to x$，放在末位。否则令 $j$ 为第一个满足 $x\to v_j$ 的下标；它大于 $1$，且 $v_{j-1}\to x$，故插入在 $v_{j-1},v_j$ 之间。单点基例显然；空图可按空序列理解。
\end{solution}

\begin{exercise}\label{mit:6042-cp17-q1}
\textbf{6.042J，Spring 2015，CP17 第 1 题。}
原卷使用的是 2006 年课程先修示例，\textbf{不是当今 MIT 的课程规定}。直接先修弧为
\[
\begin{gathered}
18.01\to6.042,18.02,18.03;\quad8.01\to8.02;\quad6.001\to6.034,6.003,6.004;\\
6.042\to6.046\to6.840;\quad18.03,8.02\to6.002;\quad
6.002\to6.003,6.004;\quad6.004\to6.033\to6.857.
\end{gathered}
\]
每门课用一学期，必须在更早学期完成先修。
\begin{enumerate}[label=(\alph*)]
\item 不限每学期门数，以每学期选全部已满足先修课程的贪心策略列出课表，证明需恰好六学期。
\item 找五门不含 18.03、能在同一学期修读的课；这样的五元组有几个？
\item 每学期只修一门，给出完整课表。
\item 每学期最多两门，最少几学期？
\item 每学期最多三门，最少几学期？
\end{enumerate}
\end{exercise}
\sourceof{CP17，物理第 1 页；编者独立解答。}
\begin{solution}
(a) 六学期的课表为
\[
\begin{array}{c|l}
1&18.01,8.01,6.001\\2&6.042,18.02,18.03,8.02,6.034\\
3&6.046,6.002\\4&6.003,6.004,6.840\\5&6.033\\6&6.857
\end{array}
\]
链 $18.01,18.03,6.002,6.004,6.033,6.857$ 的六门课必须分属六学期，证明下界。

(b) 恰有 $9$ 组：固定 $18.02,6.034,6.003$，再从 $\{6.042,6.046,6.840\}$ 和 $\{6.004,6.033,6.857\}$ 各选一门。它们两两没有间接先修，故可同学期。完整性可由五条链检查：
$\{18.01,18.03,6.002,6.003\}$，$\{8.01,8.02,6.004,6.033,6.857\}$，$\{6.001,6.034\}$，$\{6.042,6.046,6.840\}$，$\{18.02\}$。
五元反链必须每条恰选一个，故含 $18.02$，不能含 $18.01$。排除 $18.03$ 后，第一链不能选 $6.002$，因为第二链所有顶点均与它可比，故须选 $6.003$。这又排除 $8.01,8.02,6.001$，留下上列九种。可先修完该反链所有前驱，再同学期修它，故反链确实可实现。

(c) 把 (a) 各行依次读出即可，共 $15$ 学期，是拓扑序。

(d) 至少 $\lceil15/2\rceil=8$ 学期。达到下界的八行课表是
\[
\begin{array}{c|l}
1&18.01,8.01\\2&6.001,6.042\\3&18.02,18.03\\4&8.02,6.034\\
5&6.046,6.002\\6&6.003,6.004\\7&6.840,6.033\\8&6.857
\end{array}
\]
(e) 六门链仍给下界 $6$，以下达到它：第一行同 (a)，第二行 $6.042,18.03,8.02$；第三行 $18.02,6.034,6.002$；第四行 $6.046,6.003,6.004$；第五行 $6.840,6.033$；第六行 $6.857$。各行先修均在此前完成。
\end{solution}

\begin{exercise}\label{mit:6042-cp17-q2}
\textbf{6.042J，Spring 2015，CP17 第 2 题。}
两位工作人员 Lisa 和 Annie 要完成八项任务。每项只由一人连续完成，不可中断。原叙事情境是虚构的银河项目，数学数据如下：
\begin{center}
\begin{tabular}{c|l|r|l}
编号&任务&天数&直接先修\\\hline
1&设计标志与主题音乐&8&无\\2&建造舰队&18&无\\3&控制联合国&9&1\\
4&给猫接种&11&1\\5&开设咖啡连锁&10&3\\6&训练队伍&4&3,4,5\\
7&发射舰队&6&2,6\\8&击败 Microsoft&8&2,6
\end{tabular}
\end{center}
（a）给有效完成顺序；（b）用总工时给工期下界，并解释为何可能不够；（c）用关键路径给下界，并解释为何可能不够；（d）求最短工期，原题不要求证明。
\end{exercise}
\sourceof{CP17，物理第 1--2 页，图 1；编者独立解答。}
\begin{solution}
(a) $1,2,3,4,5,6,7,8$ 是有效拓扑序。

(b) 总工时为 $74$ 人天，两人给下界 $37$ 天。但先修可能使一人暂时没有可做的任务。

(c) 关键链 $1,3,5,6,8$ 的工时为 $8+9+10+4+8=39$ 天；其他最长候选链 $1,4,6,8$ 为 $31$，$2,8$ 为 $26$，故下界是 $39$。这假定关键链总能按时开始，有限工作人员还要完成链外任务。

(d) 最少 $40$ 天。Lisa 执行 $1:[0,8]$、$3:[8,17]$、$4:[17,28]$、$7:[32,38]$；Annie 执行 $2:[0,18]$、$5:[18,28]$、$6:[28,32]$、$8:[32,40]$。区间表示开始与结束时刻，任务结束后另一任务可立即开始。

先说明允许实数开工时刻也不会得到更优的非整数工期。固定任意可行方案中每位工作人员执行任务的次序，在原先修图中加入同一人相邻任务间的先后弧。所得图仍无圈：沿每条弧，后项开始时刻至少是前项开始时刻加其正工时；若有圈，把这些不等式相加就会得 $0$ 至少等于一个正数。按这个 DAG 的拓扑序，把每项的新开工时刻设为全部前驱新完工时刻的最大值，无前驱时为 $0$。整数工时使新开始、结束时刻全部为整数；沿拓扑序归纳，新的每个时刻都不晚于原方案的对应时刻。因此该方案保持人员次序和先修关系且工期不增加。特别，任何工期小于 $40$ 的方案都能变成工期至多 $39$ 的整数方案。

现在排除 $39$：若在 $39$ 完成，关键链必须无间隔地占据 $[0,8],[8,17],[17,27],[27,31],[31,39]$，链上任务即使换人，每次仅剩另一人可执行链外任务。任务 $4$ 必须在 $[8,27]$ 内占连续 $11$ 天，任务 $2$ 必须在 $33$ 之前完成，才能再做六天的任务 $7$。二者都不能与关键链同用正在工作的那个人，且不能彼此重叠；若 $2$ 在 $4$ 之前，最早 $4$ 为 $[18,29]$；若 $4$ 在 $2$ 之前，最早 $2$ 为 $[19,37]$。两者都违反期限。因此不存在 $39$ 天方案，结合整数规范化与关键链的下界，任何小于 $40$ 天的方案都不存在；上述 $40$ 天方案确实最优。
\end{solution}

\begin{exercise}\label{mit:6042-cp17-q3}
\textbf{6.042J，Spring 2015，CP17 第 3 题。}
有 $n$ 个单位时间任务，先修构成 DAG，最长链含 $t$ 个任务，$1\le t\le n$。
\begin{enumerate}[label=(\alph*)]
\item 将三种描述配到原编号术语：1 间接先修、2 拓扑序、3 链、4 反链、5 最大反链大小、6 最小反链大小、7 最长链长度、8 最短链长度。描述为：可同时做的一组任务；只有一名工作人员时的有效顺序；无限工作人员时的最短周数。
\item 在所有这样的 DAG 中，要在 $t$ 周完成所需工作人员数的最小可能值是多少？给达到它的 DAG。
\item 所需人数的最大可能值是多少？给达到它的 DAG。
\end{enumerate}
\end{exercise}
\sourceof{CP17，物理第 3 页；编者独立解答。}
\begin{solution}
(a) 依次为 $4,2,7$。同一时刻做的任务必须互无先修；单人顺序是拓扑序；无限人数可按最长前驱链长度分层，恰用 $t$ 周。

(b) 每人每周至多一项，故至少 $\lceil n/t\rceil$ 人。令 $q=\lceil n/t\rceil$，构造 $q$ 条互不相交的链，每条长至多 $t$，其中一条长 $t$。这是可能的：先分配 $t$ 项，再把剩下 $n-t\le(q-1)t$ 项分入其余 $q-1$ 条非空链（$q>1$ 时 $n-t\ge q-1$）。一人执行一条链即可。

(c) 最大可能值为 $n-t+1$。以最长前驱链长度给任务分为 $t$ 个非空层，每层为反链；某层至多 $n-(t-1)$ 项，故 $n-t+1$ 人一定够。达到它的构造是 $t-1$ 项串成链，最后一项先于剩下 $n-t+1$ 项；这些末尾任务若在第 $t$ 周完成，必须同时进行。$t=1$ 时全部任务独立，需要 $n$ 人。
\end{solution}

\begin{exercise}\label{mit:6042-cp17-q4}
\textbf{6.042J，Spring 2015，CP17 第 4 题。}
不同顶点间的弧 $a\to b$ 称为覆盖弧，若每条 $a$ 到 $b$ 的路径都使用它。令 $\mathrm{covering}(D)$ 只保留覆盖弧。
\begin{enumerate}[label=(\alph*)]
\item 原图的弧为 $12,13,14,15,16,24,26,36$，数字 $ab$ 表示 $a\to b$。找覆盖弧。
\item 证明有限 DAG 的覆盖子图与原图具有相同正长度可达关系。
\item 证明具有同一正长度可达关系的两个 DAG 有相同覆盖弧。
\item 推出覆盖子图是该可达关系下边数最少的唯一 DAG，并说明其他同关系有向图为何也不能更少。
\item 给顶点 $\{1,2\}$ 上覆盖弧相同、正可达关系不同的两图，可使用自环。
\item (i) 找顶点 $1,2,3$ 的无自环完全有向图的覆盖弧；(ii) 找定向三圈 $12,23,31$ 的覆盖弧；(iii) 比较二者的正可达关系。
\end{enumerate}
\end{exercise}
\sourceof{CP17，物理第 3--4 页，图 2；编者独立解答。}
\begin{solution}
(a) 覆盖弧为 $12,13,15,24,26,36$；$14$ 可被 $12,24$ 替代，$16$ 可被 $12,26$ 或 $13,36$ 替代。

(b) 对任意可达有序对，在有限 DAG 中选最长路径。它的每条弧必为覆盖弧：若一条弧可用至少两条弧的其他路径替代，拼接后仍无重复顶点，否则出现有向圈，于是得到更长路径，矛盾。因此覆盖子图仍连接这个有序对，逆向包含显然。

(c) 用可达关系本身刻画覆盖：$a\to b$ 是覆盖弧，当且仅当 $a$ 正可达 $b$ 且不存在不同于 $a,b$ 的 $c$ 使 $a$ 可达 $c$、$c$ 可达 $b$。若不存在这样的 $c$，可达路径只能是直接弧；若有 $c$，拼接两条路径给绕开直接弧的路径，DAG 中没有重复。这是只依赖可达关系的刻画。

(d) 任意同关系 DAG 必含全部覆盖弧，而 (b) 证明仅这些弧已足够，因此唯一最少者就是覆盖子图。具有同一正可达关系的任意有向图也必无圈：原 DAG 没有 $(v,v)$ 的正可达关系；故这里“所有有向图”实际上仍都是 DAG。

(e) 无弧图与只有自环 $1\to1$ 的图都没有不同端点的覆盖弧，但后者多了 $1$ 到自身的正可达关系。

(f)(i) 完全有向图每条弧都可由第三点绕行，故没有覆盖弧；(ii) 定向三圈的三条弧都是覆盖弧，因为从任一点出发的第一步被迫使用唯一出弧。(iii) 两图的正可达关系却都为全部九个有序对，包括各点到自身。
\end{solution}

\begin{exercise}\label{mit:6042-cp18-q3}
\textbf{6.042J，Spring 2015，CP18 第 3 题。}
设 $S=(6,4,7,9,1,2,5,3,8)$，子序列保持原位置顺序。
\begin{enumerate}[label=(\alph*)]
\item 列出所有最长递增子序列和最长递减子序列。
\item 在元素集上定义 $a\prec b$ 当且仅当 $a<b$ 且 $a$ 在 $S$ 中先于 $b$。给出其偏序图，以及极大、极小元素。
\item 说明递增、递减子序列与链、反链的关系。
\item 证明任意 $n$ 个互异实数组成的序列，存在长大于 $\sqrt n$ 的递增子序列，或长至少 $\sqrt n$ 的递减子序列。
\end{enumerate}
\end{exercise}
\sourceof{CP18，物理第 1--2 页；编者独立解答。}
\begin{solution}
(a) 最长递增子序列为 $(1,2,5,8)$、$(1,2,3,8)$，长 $4$；最长递减子序列为
\[
\begin{gathered}
(6,4,1),\quad(6,4,2),\quad(6,4,3),\\
(6,5,3),\quad(7,5,3),\quad(9,5,3).
\end{gathered}
\]
它们的长度为 $3$。
完整性可用末端递推核对，九个位置的末端递增、递减长度分别为
\[
\begin{aligned}
L^+&=(1,1,2,3,1,2,3,3,4),\\
L^-&=(1,2,1,1,3,3,2,3,2).
\end{aligned}
\]
一个最长序列的倒数第二项必须满足相应不等式且长度恰小一，逐项回溯即得到上述全部列表。

(b) 用覆盖弧确定 Hasse 图：$6\to7,4\to7,4\to5,7\to9,7\to8,1\to2,2\to5,2\to3,5\to8,3\to8$。其可达关系就是 $\prec$；极小元素为 $6,4,1$，极大元素为 $9,8$。

(c) 递增子序列恰为按原位置列出的链。反链按原位置排列时，后项不能比前项大，而元素互异，故严格递减；反过来递减子序列的任何两项都不满足 $\prec$，所以是反链。

(d) 令最长递增长度为 $h$，最长递减长度为 $w$。按各位置的递增长度分成 $h$ 层，每层是递减子序列：若较早元素较小，后者长度会更大。因此每层至多 $w$ 项，$n\le hw$。若 $h>\sqrt n$ 已成立，否则 $h\le\sqrt n$ 推出 $w\ge n/h\ge\sqrt n$，完成证明。
\end{solution}
